%% Route fresh-16.  Paste-ready; amsthm style of frontier.tex.
%% Everything below is PROVED.  Machine checks are named in \Cref{rem:fs16-checks}.
%% The full file, identical to this text, is at
%%   /tmp/claude-1000/-data-ssg-proof/d1fe2115-9b72-4784-bb94-87421ac1106c/scratchpad/fresh-16/fs16.tex

\subsection{The two qualitative sets, and why they cannot be made to answer the
threshold}\label{sec:fs16}

\Cref{def:simorder} uses the two sets $Z_0=\{w^\ast=0\}$ and
$Z_1=\{w^\ast=1\}$ and computes them by ``the standard qualitative algorithm,
from memory, unchecked''. We prove both here, inside the model, with no
arithmetic at all; the proofs are short and the algorithms are the ones the
rest of this subsection needs.

\begin{lemma}[Positive reachability]\label{lem:fs16-positive}
Let $G$ be an arbitrary SSG. Let $\mathrm{Pos}(G)$ be the least fixed point of
\[
  Y\;\longmapsto\;\{t_1\}\cup\{v\in\Vmax\cup\Vavg:\ \text{some successor in }Y\}
                    \cup\{v\in\Vmin:\ \text{both successors in }Y\}.
\]
Then $\mathrm{Pos}(G)=\{v:\val(v)>0\}$, and it is computed in time $O(N^2)$
without looking at a single probability.
\end{lemma}

\begin{proof}
Write $Y_0\subseteq Y_1\subseteq\cdots$ for the iterates and
$\mathrm{rk}(v)$ for the least $i$ with $v\in Y_i$.

($\supseteq$, i.e.\ $v\in\mathrm{Pos}\Rightarrow\val(v)>0$.)
Define $\sigma$ on $\Vmax\cap\mathrm{Pos}$ by letting $\sigma(v)$ be a
successor of strictly smaller rank; such a successor exists by the definition
of the iteration. Fix any $\tau$. From $v$ with $\mathrm{rk}(v)=i$ the token
moves to a vertex of rank $<i$ with probability $\ge\tfrac12$: at a Max vertex
$\sigma$ does it with probability $1$; at an average vertex one of the two
successors has smaller rank, so probability $\ge\tfrac12$; at a Min vertex
\emph{both} successors have smaller rank, so probability $1$. Hence
$t_1$ is reached within $N$ steps with probability at least $2^{-N}>0$.

($\subseteq$.) Let $U:=V\setminus\mathrm{Pos}$. Since $\mathrm{Pos}$ is a fixed
point, every $v\in U\cap(\Vmax\cup\Vavg)$ has \emph{both} successors in $U$ and
every $v\in U\cap\Vmin$ has \emph{some} successor in $U$; and $t_1\notin U$.
Let $\tau$ pick, at every Min vertex of $U$, a successor in $U$. Under
$(\sigma,\tau)$, for any $\sigma$, the token started in $U$ never leaves $U$,
so it never reaches $t_1$ and $\val^{\tau}(v)=0$.
\end{proof}

\begin{theorem}[Almost-sure reachability]\label{thm:fs16-assure}
Let $G$ be an arbitrary SSG. For $W\subseteq V$ let $Y(W)$ be the least fixed
point of
\[
  Y\;\longmapsto\;\{t_1\}\ \cup\
  \{v\in W\cap\Vmax: \text{some successor in }Y\}\ \cup\
  \{v\in W\cap\Vmin: \text{both successors in }Y\}
\]
\[
  \cup\ \{v\in W\cap\Vavg: \text{both successors in }W\ \text{and some successor in }Y\},
\]
and let $\mathrm{One}(G)$ be the greatest fixed point of the monotone map
$W\mapsto Y(W)$, computed by $W_0:=V$, $W_{i+1}:=Y(W_i)$; the sink $t_0$
matches no clause, so it leaves at the first step. Then
$\mathrm{One}(G)=\{v:\val(v)=1\}$, and it is computed in time $O(N^3)$ without
looking at a single probability.
\end{theorem}

\begin{proof}
Termination is clear: $Y$ is monotone in $W$ and $Y(W)\subseteq W$, so the
sequence $W_i$ decreases and stabilises within $N$ steps at the greatest fixed
point $W:=\mathrm{One}(G)$.

($\subseteq$ of the claim, i.e.\ $v\in W\Rightarrow\val(v)=1$.)
Since $W=Y(W)$, every $v\in W$ has a rank $\mathrm{rk}(v)$ in the least fixed
point computing $Y(W)$. Define $\sigma(v)$, for $v\in W\cap\Vmax$, to be a
successor of strictly smaller rank; it lies in $W$. Fix $\tau$. From $v\in W$
the token stays in $W$: at Min and average vertices of $W$ both successors are
in $W$ (at Min, because both are in $Y(W)=W$; at average, by the clause), and
at Max vertices $\sigma(v)\in W$. Moreover from every $u\in W$ the rank drops
with probability at least $\tfrac12$ in one step, as in
\Cref{lem:fs16-positive}. Hence, from every $u\in W$ and under every $\tau$,
$t_1$ is reached within $N$ steps with probability at least $2^{-N}$; by the
Markov property the probability of never reaching $t_1$ is $0$, so
$\min_\tau\Pr_{\sigma,\tau}[t_1]=1$ and $\val(v)=1$.

($\supseteq$.) Put $A:=\{v:\val(v)=1\}$. By Knaster--Tarski the greatest fixed
point of $W\mapsto Y(W)$ contains every $W$ with $W\subseteq Y(W)$, so it
suffices to prove $A\subseteq Y(A)$. Suppose not, and put
$B:=A\setminus Y(A)\ne\emptyset$; note $t_1\in Y(A)$, so $t_1\notin B$. We
record what membership in $B$ forces. If $v\in B\cap\Vavg$ then
$\tfrac12(w^\ast(v^{(0)})+w^\ast(v^{(1)}))=1$ forces both successors into $A$,
so the clause ``both successors in $A$'' holds and therefore \emph{neither}
successor is in $Y(A)$; both are in $B$. If $v\in B\cap\Vmin$ then both
successors are in $A$ and they cannot both be in $Y(A)$, so some successor is
in $B$. If $v\in B\cap\Vmax$ then every successor lying in $A$ lies in $B$ (a
successor in $A\cap Y(A)$ would put $v$ in $Y(A)$), and at least one successor
lies in $A$.

Let $c:=\max\{\val(u):u\notin A\}$, with $c:=0$ if $A=V$; then $c<1$. Fix any
Max strategy $\sigma$ and let Min play as follows: while the token is in $B$,
at a Min vertex move to a successor in $B$; as soon as the token leaves $B$,
switch to a strategy optimal at the vertex reached. Starting at $v\in B$, the
token can leave $B$ only at a Max vertex, and only into a successor not in $A$,
hence into a vertex of value at most $c$. Therefore
$\val(v)\le 0\cdot\Pr[\text{stay in }B\text{ forever}]+c\cdot\Pr[\text{leave}]\le c<1$,
contradicting $v\in A$.
\end{proof}

\begin{corollary}[The qualitative sets do not see the coins]\label{cor:fs16-bias}
Let $G$ be an SSG and let $G[b]$, for $b\in(0,1)^{\Vavg}$, be the stochastic
game obtained by letting the average vertex $v$ take $v^{(1)}$ with probability
$b_v$ and $v^{(0)}$ with probability $1-b_v$. Then
$\{v:\val_{G[b]}(v)=1\}$ and $\{v:\val_{G[b]}(v)>0\}$ are independent of $b$.
The same holds if a vertex is replaced by a controlled vertex all of whose
actions have the same two-element support.
\end{corollary}

\begin{proof}
The proofs of \Cref{lem:fs16-positive,thm:fs16-assure} use only the
\emph{supports} of the transition distributions: the constant $\tfrac12$ enters
only through the inequality ``probability at least $\tfrac12$'', which becomes
``at least $\min_v\min(b_v,1-b_v)>0$'', and the two converse arguments use no
probability at all. A controlled vertex whose actions share the support
$\{e_1,e_0\}$ satisfies exactly the same clause as a chance vertex with that
support in both fixed points.
\end{proof}

\subsubsection*{Contexts, calls and the gate lemma}

\begin{definition}[Context, call, composition]\label{def:fs16-context}
A \emph{context} is a structure $H$ on a vertex set $V_H$ together with the two
sinks, in which every vertex has an ordered pair of successors and is of one of
the four kinds $\max$, $\min$, $\mathrm{avg}$, $\mathrm{call}$; for a call $c$
we write $(e_1(c),e_0(c)):=\mathrm{succ}(c)$. $H$ is \emph{admissible} if it is
stopping when every call is read as a controlled vertex, i.e.\ if the only
subset $U\subseteq V_H$ with both successors in $U$ at every average vertex of
$U$ and some successor in $U$ at every other vertex of $U$ is $U=\emptyset$.

For a stopping SSG $G$ with start $v_0$, the \emph{composition} $H\langle
G\rangle$ is the SSG obtained from $H$ by replacing each call $c$ by a private
copy $G^{c}$ of $G$: every edge into $c$ is redirected to the copy's $v_0$,
every edge of the copy into $t_1$ is redirected to $e_1(c)$, and every edge of
the copy into $t_0$ to $e_0(c)$.
\end{definition}

\begin{definition}[The $(p,q)$-gate]\label{def:fs16-gate}
For $p,q\in[0,1]$, the \emph{$(p,q)$-gate} with exits $(e_1,e_0)$ is the vertex
whose operator is
\[
  \mathcal{G}_{p,q}(A,B)\;:=\;
  \begin{cases} pA+(1-p)B, & A\ge B,\\ qA+(1-q)B, & A<B,\end{cases}
  \qquad A:=x(e_1),\ B:=x(e_0).
\]
Equivalently $\mathcal{G}_{p,q}(A,B)=\max$ of $pA+(1-p)B$ and $qA+(1-q)B$ when
$p\ge q$, and their $\min$ when $p\le q$; so a gate is a controlled vertex with
two chance actions, each of support $\{e_1,e_0\}$. It is monotone,
$1$-Lipschitz in the sup-norm, and commutes with increasing affine maps of
$(A,B)$. We write $H^{p,q}$ for the context $H$ with every call replaced by the
$(p,q)$-gate on its exits.
\end{definition}

\begin{lemma}[Gate lemma: a plugged-in game is exactly a gate]\label{lem:fs16-gate}
Let $G$ be a stopping SSG with start $v_0$, let $G^{\mathrm{rev}}$ be $G$ with
the two players' roles exchanged (sinks unchanged), and put
\[
  p:=\val_G(v_0)=\max_\sigma\min_\tau\Pr\nolimits_{\sigma,\tau}[\,t_1\,],
  \qquad
  q:=\val_{G^{\mathrm{rev}}}(v_0)=\min_\sigma\max_\tau\Pr\nolimits_{\sigma,\tau}[\,t_1\,].
\]
Let $H$ be an admissible context. Then $H\langle G\rangle$ is stopping and
\[
  \val_{H\langle G\rangle}(v)=\val_{H^{p,q}}(v)\qquad\text{for every }v\in V_H,
\]
and inside the copy at the call $c$ the values are those of $G$ with the two
sinks paying $\val_{H^{p,q}}(e_1(c))$ and $\val_{H^{p,q}}(e_0(c))$.
No inequality between $p$ and $q$ holds in general: $p=1,q=0$ for the
one-vertex game $v\to(t_1,t_0)$ with $v\in\Vmax$, $p=0,q=1$ with $v\in\Vmin$,
and $p=q$ whenever $G$ has no controlled vertex.
\end{lemma}

\begin{proof}
\emph{Payoff transfer.} Let $G[A,B]$ be $G$ with $t_1$ paying $A$ and $t_0$
paying $B$, and let $\psi(x):=B+(A-B)x$, so that the payoff of a play is
$\psi(\mathbf 1[\,t_1\,])$; this uses that $G$ is stopping, so that a play is
absorbed almost surely and $\Pr[t_0]=1-\Pr[t_1]$. If $A\ge B$ then $\psi$ is
nondecreasing and commutes with $\max$ and $\min$, so
$\val_{G[A,B]}(v_0)=\psi(p)=pA+(1-p)B$. If $A<B$ then $\psi$ is decreasing, so
$\min_\tau\psi(\cdot)=\psi(\max_\tau\cdot)$ and
$\max_\sigma\psi(\cdot)=\psi(\min_\sigma\cdot)$, whence
$\val_{G[A,B]}(v_0)=\psi(q)=qA+(1-q)B$. In both cases
$\val_{G[A,B]}(v_0)=\mathcal{G}_{p,q}(A,B)$.

\emph{Stopping.} Let $(\sigma,\tau)$ be a positional pair of $H\langle
G\rangle$. Inside a copy it restricts to a pair of $G$, so the token leaves the
copy almost surely, landing on $e_1(c)$ or $e_0(c)$. The induced process on
$V_H$ is therefore a Markov chain in which a call moves to $e_1(c)$ with some
probability $\rho_c$ and to $e_0(c)$ with $1-\rho_c$. A finite chain absorbs
almost surely from every state iff the set of states from which no sink is
reachable is empty; that set is closed, so it has both successors in it at
every average vertex and at least one at every other vertex, and admissibility
of $H$ forces it to be empty, whatever the $\rho_c$.

\emph{Values.} Let $W:=\val_{H^{p,q}}$ and define $x$ on $V(H\langle G\rangle)$
by $x:=W$ on $V_H\setminus K$, $x:=W(c)$ at the entry of the copy of $c$, and,
on the copy of $c$, $x:=\val_{G[W(e_1(c)),W(e_0(c))]}$. At a vertex of
$V_H\setminus K$ the two Shapley operators read the same numbers, so
$(Tx)(v)=(T_{H^{p,q}}W)(v)=W(v)=x(v)$. At an interior vertex of a copy the
operator is that of $G[W(e_1),W(e_0)]$, whose value vector $x$ is, so
$(Tx)(v)=x(v)$. At the entry of the copy of $c$,
$(Tx)=\val_{G[W(e_1),W(e_0)]}(v_0)=\mathcal{G}_{p,q}(W(e_1),W(e_0))=W(c)=x$ by
the payoff transfer and the definition of $H^{p,q}$. So $Tx=x$; as
$H\langle G\rangle$ is stopping, \Cref{thm:contraction} makes $x$ the value.
\end{proof}

\begin{theorem}[Two-exit contraction]\label{thm:fs16-contract}
Let $G$ be a stopping SSG and $M$ a set of non-sink vertices such that every
edge leaving $M$ lands in $\{o_1,o_0\}$ for two fixed vertices
$o_1,o_0\notin M$. Let $G_M$ be the game induced on $M$ with $o_1$ read as
$t_1$ and $o_0$ as $t_0$ (it is stopping), and for each \emph{entry}
$e$ of $M$ --- a vertex of $M$ with an in-edge from outside $M$, or the start
vertex --- put $p_e:=\val_{G_M}(e)$ and $q_e:=\val_{G_M^{\mathrm{rev}}}(e)$.
Let $G/M$ be obtained from $G$ by deleting $M$, adding one $(p_e,q_e)$-gate
$g_e$ with exits $(o_1,o_0)$ per entry $e$, and redirecting every edge that
entered $e$ from outside $M$ to $g_e$. Then
\[
  \val_{G/M}(v)=\val_G(v)\ \ (v\notin M),\qquad \val_{G/M}(g_e)=\val_G(e),
\]
and the values inside $M$ are those of $G_M$ with the sinks paying
$\val_G(o_1)$ and $\val_G(o_0)$.

In particular the \emph{response of a part with two exits is described by two
numbers per entry}: it is the piecewise-affine map $\mathcal{G}_{p_e,q_e}$,
with exactly two affine pieces, separated by the hyperplane $A=B$.
\end{theorem}

\begin{proof}
Identical to the last paragraph of \Cref{lem:fs16-gate}, run with $W:=\val_{G/M}$
and with one gate per entry: the payoff transfer inside $M$ depends only on the
pair $(\val_G(o_1),\val_G(o_0))$ and is applied at each entry separately, and
$G_M$ is stopping because a play confined to $M$ for ever would exhibit a
nonempty trap of $G$, contradicting \Cref{lem:trapchar}. $G/M$ is stopping for
the same reason as in \Cref{lem:fs16-gate}, and $Tx=x$ with $x$ assembled as
there.
\end{proof}

\begin{remark}[Where the response table explodes, exactly]\label{rem:fs16-two}
\Cref{thm:fs16-contract} is the sharp complement of \Cref{rem:fold-width}. A
part with \emph{two} exits has a response with exactly two affine pieces, and
it can therefore be \emph{stored}: two rationals per entry. A part with
\emph{three} exits --- the situation of \Cref{thm:fold}, where the exits are
$u$, $t_0$ and $t_1$ and one free parameter survives affine normalisation ---
already has $2^{D}$ pieces on $6D+2$ vertices. So the amortisation asked for
after \Cref{prop:modulator-family} is available at separators of interface
dimension zero and is closed at interface dimension one. Consequently a
stopping SSG that can be reduced to a single vertex by repeated contraction of
parts with at most two exits and boundedly many entries, the parts and the
quotients being recursively of the same shape and of geometrically decreasing
size, is solved exactly in polynomial time; the class is narrow --- with
boundedly many entries the decomposition exhibits balanced separators of
bounded size, so it lies at bounded treewidth, where \Cref{thm:qp} already
gives a quasipolynomial bound --- and the point of the statement is not the
class but the location of the boundary.
\end{remark}

\subsubsection*{An energy identity and a crossing bound}

\begin{theorem}[Energy and crossings]\label{thm:fs16-cross}
Let $\Gamma$ be a stopping game with Max, Min and chance vertices, each chance
vertex $v$ having two successors $v^{(1)},v^{(0)}$ taken with probabilities
$b_v,1-b_v$, and two sinks of payoff $1$ and $0$. Let $W:=\val_\Gamma$, let
$(\sigma,\tau)$ be an optimal pair, let $N_v$ be the expected number of visits
to $v$ from $v_0$ in the chain $\Gamma_{\sigma,\tau}$, and put
$\Delta_v:=W(v^{(1)})-W(v^{(0)})$. Let $r+1$ be the number of distinct values
in $W(V)\cup\{0,1\}$. Then
\begin{align}
  \sum_{v\ \mathrm{chance}} N_v\,b_v(1-b_v)\,\Delta_v^{2}
  &\;=\;W(v_0)\bigl(1-W(v_0)\bigr), \label{eq:fs16-energy}\\
  \sum_{v\ \mathrm{chance}} N_v\,2b_v(1-b_v)\,|\Delta_v|
  &\;\le\;\tfrac{r}{2}+1. \label{eq:fs16-cross}
\end{align}
For an SSG (all $b_v=\tfrac12$) these read
$\sum_{v\in\Vavg}N_v\Delta_v^{2}=4w^\ast(v_0)(1-w^\ast(v_0))$ and
$\sum_{v\in\Vavg}N_v|\Delta_v|\le r+2$. Both hold verbatim with $W$ replaced by
$\val_{\sigma,\tau}$ for an \emph{arbitrary} absorbing pair $(\sigma,\tau)$.
\end{theorem}

\begin{proof}
Let $X_0=v_0,X_1,\dots$ be the chain and $M_t:=W(X_t)$. Under the optimal pair,
$W(v)=W(\sigma(v))$ at Max vertices and $W(v)=W(\tau(v))$ at Min vertices, and
$W(v)=b_vW(v^{(1)})+(1-b_v)W(v^{(0)})$ at chance vertices; so $(M_t)$ is a
bounded martingale, constant along the controlled steps, and absorbed in
$\{0,1\}$ almost surely.

\eqref{eq:fs16-energy}: at a chance vertex $v$ the increment is $(1-b_v)\Delta_v$
with probability $b_v$ and $-b_v\Delta_v$ with probability $1-b_v$, so its
conditional second moment is $b_v(1-b_v)\Delta_v^{2}$. Summing over time,
$\sum_v N_vb_v(1-b_v)\Delta_v^{2}=\mathbb{E}\sum_t(M_{t+1}-M_t)^2
=\mathrm{Var}(M_\infty)=W(v_0)(1-W(v_0))$, since $M_\infty\in\{0,1\}$ with
$\mathbb{E}M_\infty=W(v_0)$.

\eqref{eq:fs16-cross}: for $\lambda\in(0,1)$ let $C(\lambda)$ be the number of
transitions of the chain that cross the level $\lambda$, i.e.\ from $M_t\le
\lambda$ to $M_{t+1}>\lambda$ or the reverse. Only chance steps can cross. At a
chance vertex $v$ with $\Delta_v>0$ we have
$W(v^{(0)})<W(v)<W(v^{(1)})$; if $\lambda\in[W(v^{(0)}),W(v))$ then exactly the
move to $v^{(0)}$ crosses, with probability $1-b_v$, and if
$\lambda\in[W(v),W(v^{(1)}))$ then exactly the move to $v^{(1)}$ crosses, with
probability $b_v$. The two intervals have lengths $b_v\Delta_v$ and
$(1-b_v)\Delta_v$, so
\[
  \int_0^1\mathbb{E}\,C(\lambda)\,d\lambda
  =\sum_v N_v\bigl(b_v\Delta_v(1-b_v)+(1-b_v)\Delta_v b_v\bigr)
  =\sum_v N_v\,2b_v(1-b_v)|\Delta_v|,
\]
the same computation applying when $\Delta_v<0$.

Now let $0=w_0<w_1<\dots<w_r=1$ list the distinct values, so the $r$ gaps
$(w_{j-1},w_j)$ have lengths $g_j$ summing to $1$. Fix $\lambda$ in the $j$-th
gap. A crossing of $\lambda$ is a single transition from $\le w_{j-1}$ to
$\ge w_j$ or the reverse, and successive crossings alternate in direction;
hence the upward crossings of $\lambda$ are exactly upcrossings of
$[w_{j-1},w_j]$ in Doob's sense, and their number $U$ satisfies
$g_j\,\mathbb{E}U\le\mathbb{E}(M_\infty-w_{j-1})^{+}-(M_0-w_{j-1})^{+}
\le\min\bigl(w_{j-1}(1-W(v_0)),\,W(v_0)(1-w_{j-1})\bigr)\le\tfrac14$.
The downward crossings $D$ satisfy $D\le U+1$ pathwise, so
$\mathbb{E}C(\lambda)\le 2\mathbb{E}U+1\le \tfrac1{2g_j}+1$ and
\[
  \int_0^1\mathbb{E}C(\lambda)\,d\lambda
  \;\le\;\sum_{j=1}^{r} g_j\Bigl(\frac1{2g_j}+1\Bigr)\;=\;\frac r2+1 .
\]
For an arbitrary absorbing pair the same proof applies with $\val_{\sigma,\tau}$,
which is harmonic for that chain by definition.
\end{proof}

\begin{corollary}\label{cor:fs16-visits}
In a stopping SSG, every average vertex satisfies
$N_v\le 4w^\ast(v_0)(1-w^\ast(v_0))\,\Delta_v^{-2}\le\Delta_v^{-2}$: an average
vertex whose two successors differ in value by $\Delta$ is visited at most
$\Delta^{-2}$ times in expectation under an optimal pair.
\end{corollary}

\begin{corollary}[Composition is $O(|V_H|)$-Lipschitz in the plugged-in
value]\label{cor:fs16-lip}
Let $H$ be an admissible context with $s:=|V_H|$ vertices and $v\in V_H$, and
for $p,q\in(0,1)$ write $V(p,q):=\val_{H^{p,q}}(v)$. Then $V$ is continuous and
piecewise rational, and wherever it is differentiable
\[
  \Bigl|\frac{\partial V}{\partial p}\Bigr|\ \le\ \frac{s+3}{2}\cdot\frac1{2p(1-p)},
  \qquad
  \Bigl|\frac{\partial V}{\partial q}\Bigr|\ \le\ \frac{s+3}{2}\cdot\frac1{2q(1-q)} .
\]
In particular $p\mapsto V(p,p)$ is $\tfrac43(s+3)$-Lipschitz on
$[\tfrac14,\tfrac34]$.
\end{corollary}

\begin{proof}
Fix $p$ at which the optimal pair and the gate actions are locally constant.
The chain is absorbing and the standard perturbation formula for absorbing
chains gives $\partial V/\partial p=\sum_{c}N_c\bigl(W(e_1(c))-W(e_0(c))\bigr)
=\sum_c N_c\Delta_c$, the sum over the calls that use the action of bias $p$.
Those are chance steps of bias $p$ in the chain, so by
\eqref{eq:fs16-cross} --- whose left-hand side runs over all chance steps and
therefore dominates the sub-sum over those calls ---
$\sum_cN_c|\Delta_c|\,2p(1-p)\le\tfrac r2+1$ with
$r\le s+1$, since $W$ takes at most $s+2$ distinct values on $V_H\cup\{t_0,t_1\}$.
$V$ is continuous because the values of a stopping game depend continuously on
the transition probabilities, and it is piecewise rational with finitely many
pieces, one per optimal profile, so it is absolutely continuous and the bound
integrates. On $[\tfrac14,\tfrac34]$, $2p(1-p)\ge\tfrac38$.
\end{proof}

\subsubsection*{Two impossibility statements for black-box constructions}

\begin{theorem}[No black-box gap amplification]\label{thm:fs16-noamp}
Let $H$ be an admissible context on $s$ vertices, let $v\in V_H$, and for
$\theta\in(0,1)$ let $C_\theta$ be any stopping SSG without controlled vertices
whose value at its start is $\theta$ (so $p=q=\theta$ in
\Cref{lem:fs16-gate}; such a game exists for every dyadic $\theta$, a binary
chain as in \Cref{thm:compare-equivalence}). Put
$V(\theta):=\val_{H\langle C_\theta\rangle}(v)=\val_{H^{\theta,\theta}}(v)$, the
second expression defining $V$ for the non-dyadic $\theta$ as well.
If for some $\delta\in(0,\tfrac14]$
\[
  V(\theta)\ \ge\ \tfrac34\ \text{ for }\theta\ge\tfrac12,
  \qquad
  V(\theta)\ \le\ \tfrac14\ \text{ for }\theta\le\tfrac12-\delta,
\]
then $s\ \ge\ \dfrac{3}{8\delta}-3$.

Consequently: no construction that plugs copies of the given game into a fixed
context and reads the answer off a threshold on the composed value can amplify
the gap. Since by \Cref{lem:denominator-sharp} the values of a stopping SSG
with $a$ average vertices differ from $\tfrac12$ by at least $2^{-(a+1)}$ when
they differ at all, such a construction needs
$s\ge 3\cdot2^{a-2}-3=2^{\Omega(a)}$ vertices, and the composed game is
exponentially large. In particular the ``play the game $K$ times and take a
majority'' construction, and every finite-automaton verdict on a stream of
independent plays --- $H$ being the automaton, one call per state --- needs
$2^{\Omega(a)}$ states.
\end{theorem}

\begin{proof}
By \Cref{lem:fs16-gate}, $V(\theta)=\val_{H^{\theta,\theta}}(v)$, so $V$ is the
function of \Cref{cor:fs16-lip} restricted to the diagonal. The hypotheses give
$|V(\tfrac12)-V(\tfrac12-\delta)|\ge\tfrac34-\tfrac14=\tfrac12$, and
$[\tfrac12-\delta,\tfrac12]\subseteq[\tfrac14,\tfrac34]$ because
$\delta\le\tfrac14$, so \Cref{cor:fs16-lip} gives
$\tfrac12\le\tfrac43(s+3)\delta$, i.e.\ $s+3\ge 3/(8\delta)$.
\end{proof}

\begin{corollary}[The gap per vertex does not improve]\label{cor:fs16-gapratio}
With $H$ and $V$ as above, $|V(\theta)-V(\theta')|\le\tfrac43(s+3)|\theta-\theta'|$
for $\theta,\theta'\in[\tfrac14,\tfrac34]$, while $H\langle C_\theta\rangle$ has
at least $s$ vertices. So a black-box composition multiplies the quantity to be
resolved by at most $O(s)$ while multiplying the size by at least $s$: the ratio
\emph{gap\,/\,size} is not improved by more than a constant factor, whatever the
context.
\end{corollary}

\begin{theorem}[No black-box quantitative-to-qualitative
reduction]\label{thm:fs16-noqual}
Let $H$ be an admissible context and $v\in V_H$. Then for all
$p,q,p',q'\in(0,1)$,
\[
  \val_{H^{p,q}}(v)=1\iff\val_{H^{p',q'}}(v)=1,
  \qquad
  \val_{H^{p,q}}(v)>0\iff\val_{H^{p',q'}}(v)>0 .
\]
Hence no context $H$ satisfies ``$\val_{H\langle G\rangle}(v)=1$ if and only if
$\val_G(v_0)\ge\tfrac12$'' over the stopping SSGs $G$ with
$\val_G(v_0)\in(0,1)$, and likewise for $\val>0$; the same holds for any
objective whose almost-sure or positive winning set depends only on the
supports of the transition distributions.
\end{theorem}

\begin{proof}
A $(p,q)$-gate is a controlled vertex whose two actions both have support
$\{e_1,e_0\}$, independently of $p,q\in(0,1)$; apply
\Cref{cor:fs16-bias}. For the consequence, take $C_\theta$ with
$\theta<\tfrac12$ and $\theta'\ge\tfrac12$: the two composed games are
$H^{\theta,\theta}$ and $H^{\theta',\theta'}$, which have the same qualitative
sets.
\end{proof}

\begin{remark}[Repeated play is only i.i.d.\ for a one-colour
$G$]\label{rem:fs16-iid}
Call a context \emph{monotone at $(p,q)$} if $\val_{H^{p,q}}(e_1(c))\ge
\val_{H^{p,q}}(e_0(c))$ at every call. For a monotone context every gate uses
the action of bias $p$, so $\val_{H\langle G\rangle}(v)$ is exactly the
probability that the finite automaton $H$ --- one state per call, the exit
$e_1$ read on the bit $1$ --- accepts a stream of independent
$\mathrm{Bernoulli}(p)$ bits. This is the intended reading of ``play $G$ again
and again and take a verdict''. It is however \emph{false} in general: for a
two-player $G$ the copy is a $(p,q)$-gate with $p\ne q$
(\Cref{lem:fs16-gate}), the bias is $p$ at the calls whose $e_1$-exit is the
better one and $q$ at the others, and the composed value is not the acceptance
probability of any fixed automaton on i.i.d.\ bits. \Cref{thm:fs16-noamp}
covers both cases, being a statement about the diagonal $p=q$ and, through
\Cref{cor:fs16-lip}, about each parameter separately.
\end{remark}

\begin{remark}[How steep an amplifier can be, measured]\label{rem:fs16-measured}
The bound of \Cref{cor:fs16-lip} is tight up to a factor of about $8$. Over all
contexts consisting of $s$ calls only --- i.e.\ all $s$-state automata --- the
exact maximum of $\partial V/\partial p$ at $p=\tfrac12$ is
$1,\ \tfrac43,\ 2,\ \tfrac{12}5$ for $s=1,2,3,4$ (exhaustive, up to relabelling
of the states), and the best value found by hill-climbing or by the
birth-and-death context below is $3,\ \tfrac{32}9,\ 4,\ \tfrac{40}9,\ 5$ at
$s=5,\dots,9$ and $7$ at $s=13$ --- linear, not exponential, in $s$; the
birth-and-death context (the gambler's-ruin counter,
$c_j\to(c_{j+1},c_{j-1})$ with $c_0=t_0$, $c_{s+1}=t_1$, started in the middle)
attains exactly $\tfrac{s+1}2$ for odd $s$ and is the maximiser at every $s$
where the search is exhaustive. On that context the exact separations are
$V(\tfrac12+\delta)-V(\tfrac12-\delta)=
\tfrac{2080}{4481}=0.46418$, $0.46263$, $0.46225$, $0.46215$ at
$(s,\delta)=(7,\tfrac1{16}),(15,\tfrac1{32}),(31,\tfrac1{64}),(63,\tfrac1{128})$
(exact rationals, the last with a $58$-digit numerator): the product $s\delta$,
not $s$ or $\delta$ alone, governs, exactly as \Cref{thm:fs16-noamp} says. To
separate $\tfrac12$ from $\tfrac12-2^{-(a+1)}$ one therefore needs
$s=\Theta(2^{a})$ calls, hence $\Theta(2^{a})$ copies of the game.
\end{remark}

\begin{remark}\label{rem:fs16-scope}
\Cref{thm:fs16-noamp,thm:fs16-noqual} are statements about \emph{black-box}
constructions: the context is fixed and the given game is used only as a
plugged-in module entered at one vertex. They are not barriers against
arbitrary polynomial-time reductions, and cannot be: a transformation
$G\mapsto G'$ with $\val_{G'}=1$ iff $\val_G(v_0)\ge\tfrac12$ exists if and
only if \textsc{SSG-Value} is in $\mathrm{P}$, trivially in both directions
(if the problem is in $\mathrm{P}$, output the one-vertex game $t_1$ or $t_0$;
conversely \Cref{thm:fs16-assure} decides the image). What the two theorems do
is close the whole family of constructions that seed the question --- repeated
play, majority votes, tournaments, duels, gambler's-ruin counters, any
finite-state verdict on a stream of plays --- by a single Lipschitz estimate,
and they identify the exact reason: a copy of $G$ is a gate
(\Cref{lem:fs16-gate}), a gate is a chance vertex with a fixed support, and both
the almost-sure and the positive winning sets are blind to the biases while the
threshold at $\tfrac12$ is not.
\end{remark}

\subsubsection*{The mixed-Min landscape is vacuous}

\begin{lemma}[Mixing at one Min vertex never helps Min]\label{lem:fs16-mixmin}
Let $G$ be a stopping SSG with $\Vmin=\{u_1,\dots,u_k\}$. For
$y\in[0,1]^{k}$ let $G^{y}$ be $G$ with each $u_j$ replaced by a chance vertex
taking $u_j^{(1)}$ with probability $y_j$, and $V(y):=\val_{G^{y}}(v_0)$. Then
for every $j$ and every $y$,
\[
  V(y)\ \ge\ \min\bigl(V(y[y_j{:=}0]),\,V(y[y_j{:=}1])\bigr),
\]
and $\min_{y\in[0,1]^k}V(y)=\val_G(v_0)$, attained at $y=\mathbf 1_{\tau}$ for
every optimal Min strategy $\tau$.
\end{lemma}

\begin{proof}
Fix $y_{-j}$. The game $\Gamma$ obtained from $G^{y}$ by restoring $u_j$ as a
Min vertex has exactly two Min strategies, so
$\val_\Gamma(v_0)=\min(V(y[y_j{:=}0]),V(y[y_j{:=}1]))$. Pinning Min to the
mixed action $y_j$ is a particular Min behavioural strategy, and
$\val^{\tau}\ge\val$ for every $\tau$. The second statement is the same
argument applied to all coordinates at once, together with positional
determinacy.
\end{proof}

\begin{proposition}[$V$ is monotone in each coordinate]\label{prop:fs16-mixmono}
With the notation of \Cref{lem:fs16-mixmin}, for each $j$ and each fixed
$y_{-j}$ the map $y_j\mapsto V(y)$ is monotone on $[0,1]$. Consequently exact
coordinate descent on $V$ moves every coordinate to an endpoint of $[0,1]$ and
coincides, from the second round on, with single-switch Min policy iteration on
$G$: the continuous relaxation of $\Vmin$ contributes nothing.
\end{proposition}

\begin{proof}
$G^y$ has no Min vertex, so $V(y)=\max_\sigma\val_\sigma(y)$. Because $G$ is
stopping, $I-P_{\sigma,y}$ is nonsingular for every $y\in[0,1]^k$: a set closed
under the chain would be a trap of $G$ in the sense of \Cref{lem:trapchar}. The
entries of $I-P_{\sigma,y}$ are affine in $y_j$ and $y_j$ occurs only in the
row of $u_j$, so by Cramer's rule $\val_\sigma$ is a ratio of two functions
affine in $y_j$ whose denominator does not vanish on $[0,1]$; a M\"obius
function without a pole in the interval is monotone. Split the strategies into
those with $\val_\sigma$ nondecreasing and those with $\val_\sigma$
nonincreasing in $y_j$; $V$ is the maximum of one nondecreasing and one
nonincreasing function, so each sublevel set $\{V\le c\}$ is an intersection of
two intervals, i.e.\ $V$ is quasiconvex: nonincreasing on $[0,\gamma]$ and
nondecreasing on $[\gamma,1]$ for some $\gamma$. By
\Cref{lem:fs16-mixmin} the minimum of $V$ over $[0,1]$ is attained at an
endpoint; if it is at $0$ then $V$ is constant on $[0,\gamma]$ and $V$ is
nondecreasing throughout, and symmetrically. Hence $V$ is monotone.
\end{proof}

\subsubsection*{The largest controlled vertex}

\begin{theorem}[Naming the top controlled vertex is the whole
problem]\label{thm:fs16-topctrl}
Let \textsc{TopC} be the problem: given a stopping SSG with
$\Vmax\cup\Vmin\ne\emptyset$, output a controlled vertex of maximum value. Then
\textsc{TopC} is polynomial-time equivalent to \Cref{prob:main}, by a reduction
making a single oracle call and adding six vertices.
\end{theorem}

\begin{proof}
One direction is trivial. For the other it suffices, by
\Cref{thm:compare-equivalence}, to reduce the comparison of two non-sink
vertices $p,q$ of a stopping $G$ (if either is a sink the comparison is
immediate). Build $H$: redirect every edge of $G$ into
$t_1$ to a fresh average vertex $h_1$ and adjoin $h_1\to(h_2,t_0)$,
$h_2\to(t_1,t_0)$; then adjoin the average vertices $z_p\to(t_1,p)$,
$z_q\to(t_1,q)$ and the Max vertices $m_p\to(z_p,t_0)$, $m_q\to(z_q,t_0)$.

$H$ is stopping and $\val_H(h_1)=\tfrac14$, so for every positional pair the
probability of reaching $t_1$ from an old vertex is exactly one quarter of what
it was: $\val_H(v)=\tfrac14\val_G(v)$ for every old $v$, and the optimal
strategies on the old part are unchanged, a positive scaling of the payoff
commuting with $\max$ and $\min$. Hence
$\val_H(z_p)=\tfrac12+\tfrac18\val_G(p)$ and
$\val_H(m_p)=\max(\val_H(z_p),0)=\tfrac12+\tfrac18\val_G(p)$, and likewise for
$q$. Every old controlled vertex has value at most $\tfrac14<\tfrac12$, and
$z_p,z_q$ are average, so the controlled vertices of maximum value are exactly
$\{m_p\}$ if $\val_G(p)>\val_G(q)$, $\{m_q\}$ if $<$, and $\{m_p,m_q\}$ if $=$.
A single oracle call therefore decides $\val_G(p)\ge\val_G(q)$ up to the
harmless tie.
\end{proof}

\begin{remark}
\textsc{TopC} joins \Cref{thm:top} on the list of target-equivalent
reformulations, and its reduction is the simpler of the two: an average vertex
cannot be capped, which is why \Cref{thm:top} needs reference and boost chains,
whereas a Max vertex over a damped copy can be placed above everything
controlled by one leak of $\tfrac14$.
\end{remark}

\subsubsection*{Attribution, and what was machine-checked}

\begin{remark}[Prior art, from memory, unchecked against the sources]\label{rem:fs16-prior}
\Cref{lem:fs16-positive,thm:fs16-assure} are the standard qualitative
algorithms for turn-based stochastic reachability games (de~Alfaro;
Chatterjee, Jurdzi\'nski and Henzinger), which \Cref{def:simorder} already
cites from memory; what is new here is only that they are proved inside this
document's model, so that the computation of $Z_0$ and $Z_1$ no longer rests on
an unchecked citation. \Cref{cor:fs16-bias} is the folklore statement that
qualitative winning depends only on the supports.
\Cref{lem:fs16-gate,thm:fs16-contract} are two-terminal, or series, reductions
in the style of Shapley's operator equivariance and of stochastic
complementation; the only content beyond \Cref{lem:payoff-transfer} is that a
\emph{two}-exit interface collapses to the pair $(p,q)$ and to a single bit,
which is a candidate rediscovery. The perturbation formula
$\partial V/\partial p=\sum_cN_c\Delta_c$ is the classical Green's-function
derivative of an absorption probability. \eqref{eq:fs16-cross} is Doob's
upcrossing inequality; using it to bound the sensitivity of a game value to the
transition probabilities is, as far as we know, not standard, but the
consequence --- a finite-state device cannot amplify a bias $\delta$ with fewer
than $\Omega(1/\delta)$ states --- is folklore in the probabilistic-automata
literature, adjacent to Rabin's isolated-cut-point theorem, and the
\emph{moral} of \Cref{thm:fs16-noamp} is Condon's original observation that the
value gap of an SSG is exponentially small and that the natural transformations
preserve it. \Cref{lem:fs16-mixmin} is positional determinacy for Min.
\end{remark}

\begin{remark}[Machine checks]\label{rem:fs16-checks}
All in exact rational arithmetic, values always by brute force over positional
pairs or by the least-fixed-point solver, never by greedy policy iteration.
\Cref{lem:fs16-positive,thm:fs16-assure}: the two attractor computations agree
with $\{\val=0\}$, $\{\val>0\}$, $\{\val=1\}$ on $32\,000$ random SSGs (two
independent runs) of which $11\,702$ are not stopping.
\Cref{cor:fs16-bias}: $2100$ (game, bias) pairs with
random rational biases, qualitative sets unchanged. \Cref{lem:fs16-gate}:
$3820$ composed games $H\langle G\rangle$ built as honest SSGs, values by
brute force, agreeing with $\val_{H^{p,q}}$ computed independently.
\Cref{thm:fs16-contract}: $5000$ contractions of random parts with at most two
outside successors, $1800$ of them with more than one entry.
\Cref{thm:fs16-cross}: \eqref{eq:fs16-energy} verified as an exact identity on
$2500$ optimal profiles and $39\,260$ arbitrary absorbing profiles;
\eqref{eq:fs16-cross} verified on $4400$ biased games, worst ratio
$\mathrm{lhs}/\mathrm{rhs}=22/75$. \Cref{cor:fs16-lip}: the derivative formula
agrees exactly with an exact secant at $185$ of $203$ contexts (the remaining
ones being points where the optimal profile changes and the derivative is
one-sided), and $|\partial V/\partial p|$ never exceeded the bound in $900$
(context, $p$) samples, the worst ratio being $1/4$.
\Cref{rem:fs16-measured}: exhaustive over all $s$-state automata for $s\le4$.
\Cref{lem:fs16-mixmin,prop:fs16-mixmono}: $2569$ coordinate lines sampled on a
grid of nine points, no interior minimum and no non-monotone line.
\Cref{thm:fs16-topctrl}: $3500$ reductions, the value identities and the
argmax both checked. Finally both statements of \Cref{thm:fs16-cross} were
re-run on this document's engineered families rather than on random games:
$L_n$ ($n\le8$, from $v_1$ and from $w_1$), $G_m$ ($m\le7$), $\mathrm{WD}(e,j,m)$
at $N=11,16,19,23,23$ and $\mathrm{CC}(L,m)$ at $N=11,18,25,32$, the identity
\eqref{eq:fs16-energy} being exact in every case.
\end{remark}

\begin{remark}[The identity does not see the hard families]\label{rem:fs16-blind}
On $\mathrm{WD}(e,j,m)$ and on $\mathrm{CC}(L,m)$ --- the two families of
\Cref{sec:wedge,thm:hybrid-lower} --- the value at the start is exactly
$\tfrac12$, so \eqref{eq:fs16-energy} reads
$\sum_{v\in\Vavg}N_v\Delta_v^{2}=1$, and the measured total variation
$\sum_{v\in\Vavg}N_v|\Delta_v|$ is exactly $1$ as well, against a bound $r+2$
that grows linearly. The whole energy of these games sits on a single average
vertex of gap $1$ visited once. \Cref{thm:fs16-cross} therefore says nothing
about what makes them hard: it constrains the geometry of the value vector, not
the difficulty of finding it, and it is recorded here because it is what makes
\Cref{cor:fs16-lip} --- and hence the two impossibility statements --- work.
\end{remark}